\documentclass[11pt]{letter}
\usepackage[margin=1in]{geometry}
\begin{document}
\begin{letter}{Editors\\Applied Energy}
\opening{Dear Editors,}

Please consider our manuscript, ``OpenDC-LCA: boundary-aware evidence
synthesis for data-center cooling under grid decarbonization,'' for
publication as a research article in \emph{Applied Energy}.

The study addresses an energy-system decision that is increasingly important
for high-density computing: whether the lifecycle advantage of a cooling
architecture transfers across electricity systems and through grid
decarbonization. We reconstruct 24 released hyperscale data-center results and
rebuild nine EPA eGRID releases to a common AR5 GWP100 basis. The analysis
shows that the 32.46\% U.S. generation-rate decline is robust to historical
eGRID characterization changes, while local cooling rankings are more
conditional: 29.85\% of state-year screens extrapolate beyond released
electricity anchors, and a transparent boundary allowance removes the sole
2023 cold-plate/single-phase reversal.

The revision also distinguishes deterministic ranking from joint robustness.
Although two-phase immersion is first in every unperturbed state-year screen,
its U.S. first-rank frequency falls from 99.3\% to 56.2\% across declared
narrow-to-wide assumption-stress envelopes. A median 5.50\% adverse correction
to impact per equivalent useful computation erases first rank. These
frequencies are explicitly presented as sensitivity diagnostics, not
probabilities.

The methodological contribution is an open, versioned framework that connects
cooling performance, hourly integration, reliability, lifecycle inventories,
provenance, and comparative-claim checks. Code and redistributable derived
results are publicly available; a DOI-bearing immutable archive is required
before submission. The work therefore bridges energy-system
modeling, data-center thermal management, lifecycle assessment, and practical
decision support.

This manuscript is original, is not under consideration elsewhere, and has
been approved by all authors. [Confirm before submission.] We believe it fits
\emph{Applied Energy}'s emphasis on energy efficiency, energy-system
decarbonization, data analytics, and translation from analysis to
implementation.

\closing{Sincerely,\\Han Hu\\Corresponding author details to be confirmed}
\end{letter}
\end{document}
